\section{Three Deployment Decisions}\label{sec:discussion}
\paragraph{Central synthesis}
A favorable result becomes useful for deployment only after binding both \emph{where protection ends} and \emph{what the cover executes}. These are separate axes: a downstream component can generate a new observable, while an executed dummy can change the task. The three cases below turn that distinction into choices about a final emitter, an action boundary and an applicable adaptation/cost result. They retain source-local positive conclusions and identify the bridge needed for a stronger one.

\subsection{Case 1: choose the final interface to protect}
\paragraph{Decision and relevant evidence}
An application designer must decide whether upstream API padding or batching also protects the user-facing connection after a content-aware frontend. NetEcho identifies the relevant receiving path: Trace A is token counts per group; Trace B estimates token/group character lengths. A downstream application can consume decoded upstream content or stream state and emit finer-grained network observations \cite{netecho2025}. Captured deployment configurations and locally implemented active-defense settings are distinct source results. Neither exact text nor original generation timestamps are automatically recovered.

\paragraph{What changes after binding}
An upstream padded link answers a different question from protection of the final emitter. NetEcho's authors already establish re-emission in their systems. The joint synthesis specifies the next measurement: capture the last target-bearing interface rather than merely refit on the upstream trace. UI Traces offers a contrasting case: modifying stored timestamps preserves action types and order, so the first question is which surviving features predict identity \cite{uitraces2026}. Refitting that view cannot certify a different downstream one.

\paragraph{When transfer is permitted}
If a final receiver obtains the same tested projection for the same evaluated population and condition, the source-local inference result applies within that scope. A verified downstream transformation using only the protected representation, with no additional target-dependent input or state, rules out \emph{fresh-content} re-emission through that path. It does not establish zero residual leakage or transfer an empirical classifier score. Extending a formal protection bound would additionally require that bound to apply to the verified transformation. Otherwise, the new final-hop implementation needs its own evaluation. This conditional permission concerns a specified path, not a measured new deployment.

\subsection{Case 2: choose whether cover may execute}
\paragraph{Decision and relevant evidence}
A designer considering dummy tool calls or GUI actions needs task/policy evidence as well as communication costs. Dangol's live OHTTP/A2A binding, with client, relay, gateway and agent co-located on one loopback host, reports workflow-classification accuracy 0.30 at about 13$\times$ bytes and 11$\times$ latency. Its transformed-capture frontier instead reports six-class accuracy 0.17 at 16.2$\times$ bytes \cite{dangol2026workflow}. AHB executes mobile humanization: Trip Planning success is 0.75 Raw versus 0.15 for the combined Online+Fake condition, with unintended GUI transitions discussed by the authors \cite{zhu2026turing}. That contrast is not an isolated causal effect of Fake alone.

\paragraph{What changes after binding}
Transport cover and semantic action cover have different outcome objects. Dangol's measured bytes/latency price its protocol condition; they do not establish that dummy GUI actions preserve an objective. Conversely, AHB's task loss is not a cost estimate for discarded transport dummies. AHB studies human--agent distinction, so this example supports a state-effect comparison rather than automatically demonstrating private-content leakage. The authors establish the local costs and humanization trade-offs; the synthesis identifies which deployment inference cannot cross the execution boundary.

\paragraph{When transfer is permitted}
Verified discard before semantic dispatch, isolated non-effecting shadow execution, or demonstrated goal/policy-neutral actions excludes the direct fake-touch mechanism within that scope. Resource use, delay and timeout effects still require evidence. Source-local outcomes also remain usable: AHB's online History+Long configuration without Fake reports task success 0.7143, but XGBoost detection remains 0.9905. This does not prove noninferiority to Raw, effective indistinguishability or a full cost frontier. It does prevent generalizing the 0.15 outcome to every humanization configuration.

\subsection{Case 3: choose a valid adaptation result and price}
\paragraph{Decision and relevant evidence}
A service owner choosing cache protection under a workload budget needs both a valid protected-condition attack estimate and a cost for the actual matching mechanism. An adaptation result from another receiving view answers only its own subproblem. UI Traces compares clean- and delay-trained classifiers on the same delayed archive view. That permits the positive source-local conclusion that adaptation recovers substantial prediction under the tested transformation (D2) \cite{uitraces2026}. SIGLEAK's useful procedural inference under a common SkillGuard5 prompt is a different result, not an installation effect or a paired adaptation contrast \cite{geng2026sigleak}.

Early Bird's $K$-specific hit/miss calibration and subsequent KV-oracle simulation are distinct from its GPTCache de-identification prototype's 6ms/4\% overhead \cite{earlybird2025}. KVGov measures a 1.3/1.6$\mu$s median/p99 HMAC increment, while workload-dependent lost cache reuse remains unmeasured; key-equality isolation does not establish isolation of a semantic nearest-neighbor index \cite{addagada2026kvgov}. SpliceDefense likewise separates CTBF's tested signal effect from QCP's padding cost \cite{spliceleak2026}. For full-prompt recovery, Early Bird Table~VI also distinguishes queries per recovered token from recovery coverage and all-token cost; the authors already explain this survivor-conditioned denominator.

\paragraph{What changes, and when transfer is permitted}
Allow the UI adaptation contrast for its matched target, receiver, transformation and defended test units. Allow each component cost for the measured component and mechanism. A new receiver, workload, combined mechanism or matching primitive needs a corresponding bridge. The original authors supply the measurements; our added output is the explicit choice of which result may inform which decision. Comparing the smallest reported milliseconds or microseconds would answer no common deployment question.

\subsection{Evidence strength and the resulting decision rule}
Table~\ref{tab:synthesis-decisions} distinguishes direct support, extensions and background; roles apply to a particular claim, not the quality of a paper. Preservation has two different instances (UI actions and SIGLEAK procedures); NetEcho alone directly supports downstream re-emission here, and AHB alone supplies the measured mobile state-change instance. Dangol supplies a contrasting transport result, not a replication of AHB. The 60-family total does not multiply those direct-source counts. Source deletion and unresolved-field sensitivity are reported in Appendix~\ref{app:claim-sensitivity}.
\begin{table*}[t]\centering\footnotesize
\caption{Three deployment judgments and their evidence roles. Direct = the named result tests that branch; extension = a different mechanism, mode or component; background = a general principle. Conditional permissions are bridge obligations, not new measurements.}\label{tab:synthesis-decisions}
\begin{tabularx}{\textwidth}{@{}p{.16\textwidth}p{.27\textwidth}YY@{}}\toprule
Decision & Direct support / extension / background & Judgment after binding & Positive license and remaining bridge\\\midrule
Protect the final stream & Direct: NetEcho IV-C--E/Fig.5 (one re-emission source). Extension: UI Traces 6.1/Fig.5, surviving actions. Background: compound-system composition \cite{banerjee2024compound}. & An upstream-view result does not certify a new content-aware emitter. & Keep the result for the tested view. Verify the same projection, or evaluate the final hop; protected-only postprocessing rules out fresh-content emission, not all leakage.\\
Implement cover as actions & Direct: AHB 5.2.3/Table1/E.3 (one mobile state-change source). Extension: Dangol \S9.7, OHTTP paragraph after Table~5, transport costs. Background: traffic-shaping trade-offs \cite{cai2014systematic}. & Bytes/latency cannot establish objective/policy-preserving dummy actions. & Keep measured protocol costs and configuration-specific task points. Verify discard/shadow semantics or measure executed goal/policy and resource effects.\\
Choose adaptation and cost & Direct: UI Traces 6.1/Fig.5, matched delayed-view adaptation. Extension: Early Bird V-A/B, KVGov 5.9/7, SpliceDefense VII-C component/mode evidence. Background: protection boundaries \cite{forough2026secrets}. & Permit the source-local adaptation effect and component price; withhold a cross-component deployment ranking. & Bind training access and defended test units; match cache-resolution primitive and workload before transferring mechanism or cost.\\\bottomrule
\end{tabularx}\end{table*}


The practical rule is to locate the final receiver, follow the intervention to its execution boundary, then bind inference, utility and cost to that condition. It permits a scoped positive judgment where evidence matches and names the exact observation that would change an unresolved judgment. The cases are proposed deployment decisions, not documented mistakes or contradictions attributed to the source authors.
